% Evans 11.1, printed pp.571--582; English PDF pp.588--599.
% Chinese section start: printed p.404 / scan p.422.
\section{Introduction}
\label{sec:11-1}
\subsection{Conserved variables and the weak formulation}
Let $\Omega\subset\R^m$ be an open state domain and let
$F\in C^3(\Omega;\R^m)$. We study
\begin{equation}\label{eq:11-system}
 u_t+F(u)_x=0,\qquad u(x,0)=g(x),\quad x\in\R.
\end{equation}
For a smooth solution, integration over $[a,b]$ gives
$\frac d{dt}\int_a^b u\,dx=F(u(a,t))-F(u(b,t))$.
This is a vector balance: each component is a conserved density.
In several space dimensions the flux is an $m\times n$ matrix
and its divergence is taken row by row.
\PDEAnnotation{守恒变量的选择}{浅水取 $u=(h,q)$、$q=hv$。
速度 $v$ 是原始变量而不是动量密度。光滑时可作变量替换，
跨激波时必须回到原守恒律的跳跃条件。}

\begin{definition}[Integral solution; source \S11.1.1]
A bounded measurable $u$ with essential range in a compact subset of
$\Omega$ is an integral solution if, for every
$\varphi\in C_c^\infty(\R\times[0,\infty);\R^m)$,
\begin{equation}\label{eq:11-weak}
 \int_0^\infty\!\int_{\R}
 [u\cdot\varphi_t+F(u)\cdot\varphi_x]\,dx\,dt
 +\int_{\R} g\cdot\varphi(x,0)\,dx=0.
\end{equation}
A strong $L^1_{\rm loc}$ initial trace, when required, is an
additional property beyond this distributional initial condition.
\end{definition}

\begin{theorem}[Jump condition; source \S11.1.1, equation (20)]
Suppose $u$ is piecewise $C^1$ across $x=s(t)$, $s\in C^1$, with
one-sided traces, and solves the PDE on each side. Put
$[u]=u_r-u_l$, $[F]=F(u_r)-F(u_l)$. The interior weak equation
holds if and only if
\begin{equation}\label{eq:11-rh}
 [F]=s'(t)[u].
\end{equation}
Together with the correct initial trace this implies \eqref{eq:11-weak}.
\end{theorem}
\begin{expandedproof}[editorial distribution calculation]
Integrate separately on $x<s(t)$ and $x>s(t)$ on a compact time strip.
The singular part of $u_t$ is $-s'[u]\delta_{x=s(t)}$; that of
$F(u)_x$ is $[F]\delta_{x=s(t)}$. All regular terms vanish by the
classical equations. Arbitrary vector tests on the curve force every
component of $[F]-s'[u]$ to vanish. Integrating from $t=\delta$
and sending $\delta\downarrow0$ supplies the initial term if the
trace converges in $L^1_{\rm loc}$. Values assigned on the curve
itself do not change any integral.
\end{expandedproof}
\PDEAnnotation{同一个激波速度}{若某分量 $[u^i]=0$，则须有
$[F^i]=0$；不能在该分量中作除法。所有非零分量给出的
$[F^i]/[u^i]$ 必须相同，故系统中并非任意两状态都能接激波。}

\subsection{Physical models and characteristic speeds}
For the elastic $p$-system use strain $a$ and velocity $b$:
\[
 a_t-b_x=0,\qquad b_t-p(a)_x=0,\qquad
 A=DF=\begin{pmatrix}0&-1\\-p'(a)&0\end{pmatrix}.
\]
This follows from $a=y_x$, $b=y_t$ for
$y_{tt}-(p(y_x))_x=0$. If $p'>0$, the speeds are
$\lambda_1=-c$, $\lambda_2=c$, $c=\sqrt{p'}$, with right vectors
$(1,c)$ and $(1,-c)$.
\PDEAnnotation{本章的应力符号}{这里 $p(a)$ 随应变增加，故
$p'>0$。气体的 Lagrange 比容形式常采用另一种符号约定，
不能把那里的压强导数符号直接搬来。}
For shallow water with gravity scaled to one,
\begin{equation}\label{eq:11-sw}
 h_t+q_x=0,\qquad q_t+(q^2/h+h^2/2)_x=0,\qquad h>0.
\end{equation}
The speeds are $v\mp\sqrt h$, $v=q/h$. The strict gap
$2\sqrt h$ closes at vacuum.
For Euler, the conserved variables are $(\rho,\rho v,\rho E)$,
$E=e+v^2/2$, and the flux is
$(\rho v,\rho v^2+p,(\rho E+p)v)$. On $\rho>0$, smooth solutions obey
\[
 \rho_t+v\rho_x+\rho v_x=0,\quad
 v_t+vv_x+\rho^{-1}p_x=0,\quad
 e_t+ve_x+(p/\rho)v_x=0.
\]
Subtracting the kinetic-energy balance from total energy gives the
last equation. Their matrix is $vI+C$, where
\[
 C=\begin{pmatrix}0&\rho&0\\p_\rho/\rho&0&p_e/\rho\\
 0&p/\rho&0\end{pmatrix},\qquad c^2=p_\rho+pp_e/\rho^2.
\]
Its characteristic polynomial is $\det(\mu I-C)=\mu(\mu^2-c^2)$.
Thus $c^2>0$ gives speeds $v-c,v,v+c$. Evans's assumptions
$p,p_\rho,p_e>0$ are sufficient; division by $\rho$ requires positivity.

\subsection{Strict hyperbolicity and smooth eigenframes}
\begin{definition}[Source \S11.1.2]
The system $u_t+B(u)u_x=0$ is strictly hyperbolic if $B(z)$ has
real distinct eigenvalues $\lambda_1<\cdots<\lambda_m$.
Write $Br_k=\lambda_k r_k$, $l_kB=\lambda_k l_k$ and choose
$l_j\cdot r_k=\delta_{jk}$. Right eigenvectors are unit when discussing
arc length; the dual left eigenvectors need not also be unit.
\end{definition}
For $j\ne k$, $(\lambda_j-\lambda_k)l_j\cdot r_k=0$.
The right eigenvectors form a basis, so $l_k\cdot r_k\ne0$ before
normalization. A smooth inviscid traveling wave $V(x-\sigma t)$
satisfies $(B(V)-\sigma I)V'=0$. Where $V'\ne0$, its speed must be
an eigenvalue along its path. A viscous profile instead uses the
scaled variable $(x-\sigma t)/\varepsilon$.

\begin{theorem}[Coordinate invariance; source \S11.1.2, Theorem 1]
Let $H$ be a $C^2$ diffeomorphism of state domains. For smooth $u$,
$\widetilde u=H(u)$ satisfies a system with matrix
\[
 \widetilde B(H(z))=DH(z)B(z)DH(z)^{-1}.
\]
Strict hyperbolicity and genuine nonlinearity are preserved.
Conservation form need not be preserved.
\end{theorem}
\begin{expandedproof}[chain rule and similarity]
Multiply the PDE by $DH(u)$ and use $\widetilde u_x=DH(u)u_x$.
Similarity preserves eigenvalues. Right vectors become $DH r_k$;
left vectors become $l_k DH^{-1}$. Moreover
$D_{\widetilde z}\widetilde\lambda_k\cdot(DH r_k)
=D_z\lambda_k\cdot r_k$. Rescaling a vector by a nonzero factor
preserves whether this expression is zero. This proof uses a
classical chain rule. Through a jump, the corresponding product of
a discontinuous coefficient and a measure is not justified.
\end{expandedproof}

\PDEAnnotation{局部与全局特征向量}{隐函数定理先给局部选择。
一般有孔状态域的特征直线丛可能无法连续定向；本章波曲线和
Riemann 拼接只需要一个小邻域。}
\begin{theorem}[Eigenframe dependence; source \S11.1.2, Theorem 2]
A smooth strictly hyperbolic matrix field has smooth ordered
 eigenvalues and local smooth eigenframes. On state domain $\R^m$,
unit eigenvectors can be chosen globally and smoothly. On a general
open state domain the eigenframe assertion made here is local.
\end{theorem}
\begin{expandedproof}[including the global orientation step]
At a fixed state take $Br_0=\lambda_0r_0$, $|r_0|=1$.
Apply the implicit function theorem to
$(B(z)r-\lambda r,(|r|^2-1)/2)$. The kernel of its derivative in
$(r,\lambda)$ satisfies
$(B-\lambda_0I)h-\mu r_0=0$, $r_0\cdot h=0$.
Multiplication by $l_0$ gives $\mu=0$; simplicity gives $h=cr_0$,
then $h=0$. The square derivative is invertible. Ordered eigenvalues
agree on overlaps, so they are global.

For global eigenvectors on $\R^m$, form the smooth projection
\[
 P_k(z)=\prod_{j\ne k}
 \frac{B(z)-\lambda_j(z)I}{\lambda_k(z)-\lambda_j(z)}.
\]
Put $P(s)=P_k(sz)$ and solve
$R_s=[P_s,P]R$, $R(0)=r_k(0)$, $0\le s\le1$.
Differentiating $P^2=P$ gives $PP_sP=0$ and
$[[P_s,P],P]=P_s$. If $U_s=[P_s,P]U$, $U(0)=I$, then
$P(s)U(s)=U(s)P(0)$: subtract the two sides, differentiate, and use
this commutator identity and uniqueness for a homogeneous linear ODE.
Hence $R$ remains in the range of $P$, and never vanishes since
$U$ is invertible. Smooth dependence on $z$ and normalization at
$s=1$ give a global unit vector. Apply the same construction to the
transpose and rescale the left vectors to be dual. This radial
construction uses the contractibility of $\R^m$.
\end{expandedproof}
